\documentclass[12pt]{article}
\usepackage[utf8]{inputenc}
\usepackage[T1]{fontenc}
\usepackage{hyperref}
\title{Taungyo Language Resources for K-12 Curriculum Development}
\author{Compiled by Claude Sonnet 5 (claude-sonnet-5) for Andr\'as Kornai}
\date{}
\begin{document}
\maketitle
\section{Introduction}
Taungyo (ISO 639-3: tco; Glottocode taun1248) is a Burmish language of Myanmar. Glottolog places it in Southern Burmish, and the EGIDS level is 6a (vigorous). Joshua Project estimates about 46,000 speakers in Kalaw and Pinlaung townships, southern Shan State. Other names are Taru, Taung-yo, Toru and Twinreu. Some sources treat it as a non-standard variety of Burmese. Wikitongues lists it as unwritten.
\section{Primary Educational Resources}
I found no primers, textbooks or school programmes. Burmese materials would be the nearest relevant educational resources, but I did not search for them.
\section{Linguistic and Reference Materials}
\begin{itemize}
\item Glottolog languoid page for Taungyo. \url{https://glottolog.org/resource/languoid/id/taun1248.bigmap.html}
\item Glottolog entry by ISO code. \url{https://glottolog.org/resource/languoid/iso/tco}
\item Wikitongues page for Taungyo. \url{https://wikitongues.org/languages/tco/}
\end{itemize}
\section{Religious and Community Resources}
\begin{itemize}
\item Scripture Earth lists Bible reading resources and audio recordings for tco. \url{https://www.scriptureearth.org/00eng.php?iso=tco}
\item Joshua Project, Taungyo in Myanmar. \url{https://joshuaproject.net/people_groups/15290/BM}
\item Asia Harvest profile PDF, Taungyo. \url{https://operation.asiaharvest.org/MYANMAR-Taungyo.pdf}
\end{itemize}
\section{Digital and Audio Resources}
Scripture Earth (above) is the only source found that offers audio.
\section{Curriculum Development Context}
The language is vigorous and has tens of thousands of speakers, but I found no standard orthography and no school use. I found no dictionary, grammar or teaching materials. Any mother-tongue curriculum would first need an agreed writing system, which should be settled with the community. The Scripture Earth materials show that some written and audio material exists. I did not examine them, so I cannot say what script or orthography they use.
\end{document}